% ch-a4-18.tex —— 第 4 篇（Emerton--Gee--Hellmann）续片，底本 PDF 384--395 = 印刷页 371--382
\ysrc{371}
于是我们可以考虑如下方块
\begin{center}
\begin{tikzcd}
\widehat{D(\mathcal A)} \arrow[r, "\widehat F"] & \widehat{D_{\mathcal A}(\mathcal B)} \\
D(\mathcal A) \arrow[r, "F"'] \arrow[u, "\sim"] & D_{\mathcal A}(\mathcal B) \arrow[u, "\sim"]
\end{tikzcd}
\end{center}
其中帽子表示左完备化（left completion）；竖直箭头是从所指范畴到其左完备化的典范函子，正如我们刚刚注意到的，它们都是等价；下水平箭头是函子 $F$，上水平箭头是 $F$ 在左完备化上诱导的函子 $\widehat F$。我们断言此图交换（差一个自然同构）。由于 $\widehat F$ 是等价（因为对每个 $n$，$F$ 都限制为等价 $D(\mathcal A)^{\ge n}\to D_{\mathcal A}(\mathcal B)^{\ge n}$），我们断定下水平箭头是等价，如所欲证。

为看出所断言的交换性，注意：若 $X$ 是 $D(\mathcal A)$ 的对象，它在 $\widehat{D(\mathcal A)}$ 中对应元素 $\lim\tau^{\ge n}X$，则有自然同构
\[
\widehat F(\lim\tau^{\ge n}X):=\lim F(\tau^{\ge n}X)\isom\lim\tau^{\ge n}F(X),
\]
其中用到 $F$ 的 $t$-正合性；而 $\lim\tau^{\ge n}F(X)$ 恰是 $D_{\mathcal A}(\mathcal B)$ 的对象 $F(X)$ 在 $\widehat{D_{\mathcal A}(\mathcal B)}$ 中的像。

\end{proof}

% 说明：本片开头直至上一段，是 ch-a4-17 所开的「命题 A.7.3 证明」的续文（底本在此段末才有证明结束的 □）；
% 命题 A.7.3 的证明在 ch-a4-17 中开启，其 \end{proof} 已移到本片中该段之后（见上）。
% 拼装已完成，无需再手工移动。

\begin{rem}{A.7.4}\label{thm:a4-18-1}
在命题 A.7.3(2) 的语境中，上述证明表明 $D(\mathcal A)$ 与 $D(\mathcal B)$ 都是左完备的，并且 $F$ 从 $D^+(\mathcal A)$ 到 $D(\mathcal A)$ 的延拓可按公式 $F(X):=\lim F(\tau_{\ge n}X)$ 计算。
\end{rem}

\subsection*{A.8.\quad 半正交分解}
本节概述稳定 $\infty$-范畴的半正交分解（semiorthogonal decomposition）的基本概念。所考虑的稳定 $\infty$-范畴一律余完备。此外，当我们说「余完备子范畴」时，一律指最强的意义，即所论子范畴在所在范畴中对取余极限封闭。

我们先证明关于余完备稳定 $\infty$-范畴中生成元的一些一般事实。

\begin{defn}{A.8.1}\label{thm:a4-18-2}
设 $\mathcal C$ 是稳定 $\infty$-范畴，$X$ 是 $\mathcal C$ 的一个对象集合。定义左正交补与右正交补
\[
{}^{\perp}X:=\{y\in\mathcal C\mid \operatorname{RHom}(y,x)=0\ \text{对一切}\ x\in X\},
\]
\[
X^{\perp}:=\{y\in\mathcal C\mid \operatorname{RHom}(x,y)=0\ \text{对一切}\ x\in X\}.
\]
\end{defn}

\begin{lem}{A.8.2}\label{thm:a4-18-3}
若 $\mathcal C$ 是余完备稳定 $\infty$-范畴，$X$ 是 $\mathcal C$ 的一个对象集合，则 ${}^{\perp}X$ 是 $\mathcal C$ 的余完备稳定子范畴。
\end{lem}
\begin{proof}
${}^{\perp}X$ 的余完备性直接由
\[
\operatorname{RHom}\bigl(\operatorname*{colim}_i y_i,x\bigr)=\lim_i\operatorname{RHom}(y_i,x)
\]
得到，故只需验证稳定性。由 [Lur17, Lem. 1.1.3.3]，只需验证 $\mathcal C$ 在取余纤维与平移下封闭。对余纤维的断言直接由定义得到；\ysrc{372}对平移的断言则来自等式
\[
\operatorname{RHom}(y[1],x)=\Omega\operatorname{RHom}(y,x),
\]
它表明 $y\in{}^{\perp}X$ 当且仅当 $y[1]\in{}^{\perp}X$。
\end{proof}

\begin{lem}{A.8.3}\label{thm:a4-18-4}
若 $\mathcal C$ 是余完备稳定 $\infty$-范畴，$X$ 是 $\mathcal C$ 中一个由紧对象组成的集合，则 $X^{\perp}$ 是 $\mathcal C$ 的余完备稳定子范畴。
\end{lem}
\begin{proof}
每个 $x\in X$ 都紧，故由注 A.4.2 有
\[
\operatorname{RHom}(x,\operatorname*{colim}_i y_i)=\operatorname*{colim}_i\operatorname{RHom}(x,y_i),
\]
可见 $X^{\perp}$ 余完备。由定义立即可见 $X^{\perp}$ 在取纤维下封闭；结论的其余部分与引理 A.8.2 的证明相同（改用 [Lur17, Lem. 1.1.3.3] 关于纤维而非余纤维的变体，该变体可由前引结果过渡到反范畴得到）。
\end{proof}

我们将用引理 A.8.2 来定义由 $X$ 生成的 $\mathcal C$ 的子范畴 $\langle X\rangle$。在此之前，先考虑如下更一般的情形会比较方便。设 $F:\mathcal C\to\mathcal D$ 是余完备稳定 $\infty$-范畴之间的全忠实连续函子。由引理 A.4.6，函子 $F$ 有右伴随 $G:\mathcal D\to\mathcal C$。以 $\mathcal A\subseteq\mathcal D$ 记 $G$ 的核，以 $\mathcal B\subseteq\mathcal D$ 记 $F$ 的本质像。二者都是 $\mathcal D$ 的稳定子 $\infty$-范畴。

\begin{lem}{A.8.4}\label{thm:a4-18-5}
在上述情形中，有如下结论。
\begin{enumerate}[label=(\arabic*),leftmargin=2.4em,itemsep=0.15em]
\item $F$ 诱导等价 $\mathcal C\isom\mathcal B$。
\item $G$ 诱导等价 $\mathcal D/\mathcal A\isom\mathcal C$。
\item 复合 $\mathcal B\subseteq\mathcal D\to\mathcal D/\mathcal A$ 是等价。
\item $\mathcal B$ 在 $\mathcal D$ 中对取余极限封闭。
\item $\mathcal A=\mathcal B^{\perp}$。
\item $\mathcal B={}^{\perp}\mathcal A$。
\end{enumerate}
\end{lem}
\begin{proof}
(1) 直接由 $\mathcal B$ 作为 $F$ 的本质像的定义以及 $F$ 全忠实的假设得到。为证 (2)，先注意：按 $\mathcal A$ 作为 $G$ 的核的定义，函子 $G$ 诱导出函子 $\mathcal D/\mathcal A\to\mathcal C$。由于伴随的单位 $\mathrm{id}_{\mathcal C}\isom GF$ 是等价（由 $F$ 全忠实），而余单位在模去 $\mathcal A$ 的意义下是等价（即给出等价 $GFG\isom G$），我们断定这个诱导出的函子是等价。(3) 只是用 (1) 与 (2) 对事实 $\mathrm{id}_{\mathcal C}\isom GF$ 的重新表述。(4) 由 $F$ 连续这一事实得到。

为证 (5)，注意 $Y\in\mathcal A$ 当且仅当 $G(Y)=0$，当且仅当对一切 $X\in\mathcal C$ 有 $\operatorname{RHom}(X,G(Y))=0$，当且仅当对一切 $X\in\mathcal C$ 有 $\operatorname{RHom}(F(X),Y)=0$（由 $F$ 与 $G$ 的伴随性），当且仅当 $Y\in\mathcal B^{\perp}$。为证 (6)，先注意由 (5) 可得包含 $\mathcal B\subseteq{}^{\perp}\mathcal A$。现设 $Y\in\mathcal D$，令 $Z$ 是伴随的余单位 $FG(Y)\to Y$ 的余纤维。由于 $GF\isom\mathrm{id}_{\mathcal C}$，态射 $GFG(Y)\to G(Y)$ 是等价，故 $Z\in\mathcal A$。特别地，若 $Y\in{}^{\perp}\mathcal A$，则由 $FG(Y)\in\mathcal B\subseteq{}^{\perp}\mathcal A$ 可知也必须有 $Z\in{}^{\perp}\mathcal A$；又因 $Z\in\mathcal A$，便知 $\mathrm{id}_Z=0$，从而 $Z=0$，如所欲证。
\end{proof}

\begin{defn}{A.8.5}\label{thm:a4-18-6}
设 $\mathcal C$ 是余完备稳定 $\infty$-范畴，$X$ 是 $\mathcal C$ 的一个对象集合。定义由 $X$ 生成的子范畴为 $\langle X\rangle:={}^{\perp}(X^{\perp})$。若 $\langle X\rangle=\mathcal C$，则称 $X$ 是 $\mathcal C$ 的一个生成元集。
\end{defn}

\ysrc{373}这个定义看起来可能相当奇怪，但注意由引理 A.8.2，$\langle X\rangle$ 是 $\mathcal C$ 的余完备稳定子范畴。下面的引理表明它是最小的、包含 $X$ 的余完备稳定子范畴，这为上述定义提供了依据。

\begin{lem}{A.8.6}\label{thm:a4-18-7}
设 $\mathcal C$ 是余完备稳定 $\infty$-范畴，$X$ 是 $\mathcal C$ 的一个对象集合，$\mathcal E$ 是 $\mathcal C$ 的余完备稳定子范畴且包含 $X$ 中的每个对象。则 $\mathcal E$ 包含 $\langle X\rangle$。
\end{lem}
\begin{proof}
取引理 A.8.4 中的 $F$ 为包含函子 $\mathcal E\subseteq\mathcal C$，可得 $\mathcal E={}^{\perp}(\mathcal E^{\perp})$。由于（按定义）$\langle X\rangle={}^{\perp}(X^{\perp})$，只需证明 ${}^{\perp}(X^{\perp})\subseteq{}^{\perp}(\mathcal E^{\perp})$。按定义，只需证明 $X^{\perp}\supseteq\mathcal E^{\perp}$，而由假设 $\mathcal E$ 包含 $X$，这立刻可得。
\end{proof}

\begin{cor}{A.8.7}\label{thm:a4-18-8}
设 $\mathcal D$ 是余完备稳定 $\infty$-范畴，$X$ 是 $\mathcal D$ 的一个对象集合。则以下断言等价。
\begin{enumerate}[label=(\arabic*),leftmargin=2.4em,itemsep=0.15em]
\item $X$ 是 $\mathcal D$ 的一个生成元集。
\item 包含 $X$ 中每个对象的 $\mathcal D$ 的余完备稳定子范畴只有 $\mathcal D$ 自身。
\item 若 $\mathcal D$ 的对象 $y$ 对每个 $x\in X$ 都满足 $\operatorname{RHom}(x,y)=0$，则 $y=0$。
\end{enumerate}
\end{cor}

现在回忆半正交分解的含义（我们的定义可能与文献中的其他定义略有出入：我们允许分解是无限的）。设 $\mathcal C$ 是余完备稳定 $\infty$-范畴，$\mathcal A_1,\mathcal A_2,\dots$ 是（有限或无限个）余完备稳定子范畴。

\begin{defn}{A.8.8}\label{thm:a4-18-9}
称 $\mathcal A_1,\mathcal A_2,\dots$ 是 $\mathcal C$ 的一个半正交分解，如果
\begin{enumerate}[label=(\arabic*),leftmargin=2.4em,itemsep=0.15em]
\item 子范畴 $\mathcal A_1,\mathcal A_2,\dots$ 生成 $\mathcal C$，且
\item 当 $x\in\mathcal A_i$、$y\in\mathcal A_j$ 且 $i<j$ 时，$\operatorname{RHom}(x,y)=0$。
\end{enumerate}
\end{defn}

\begin{rem}{A.8.9}\label{thm:a4-18-10}
按定义，定义 A.8.8 中的条件 (2) 等价于要求：若 $i<j$，则 $\mathcal A_i\subseteq{}^{\perp}\mathcal A_j$；也等价于要求：$\mathcal A_j\subseteq\mathcal A_i^{\perp}$。
\end{rem}

我们也可以用生成元来定义半正交分解，如下面的定义与引理所示。

\begin{defn}{A.8.10}\label{thm:a4-18-11}
设 $\mathcal C$ 是余完备稳定 $\infty$-范畴。称 $\mathcal C$ 中一个（可能无限的）对象序列 $x_1,x_2,\dots$ 为弱例外序列（weakly exceptional），如果：
\begin{enumerate}[label=(\arabic*),leftmargin=2.4em,itemsep=0.15em]
\item 各个 $x_i$ 都是 $\mathcal C$ 的紧对象。
\item 各个 $x_i$ 生成 $\mathcal C$。
\item 当 $i<j$ 时，$\operatorname{RHom}(x_i,x_j)=0$。
\end{enumerate}
\end{defn}

\begin{lem}{A.8.11}\label{thm:a4-18-12}
设 $\mathcal C$ 是余完备稳定 $\infty$-范畴，$x_1,x_2,\dots$ 是 $\mathcal C$ 中一个弱例外对象序列，并令 $\mathcal A_i:=\langle x_i\rangle$ 是由 $x_i$ 生成的子范畴。则 $\mathcal A_1,\mathcal A_2,\dots$ 是 $\mathcal C$ 的一个半正交分解。
\end{lem}
\begin{proof}
按假设，序列 $x_1,x_2,\dots$ 生成 $\mathcal C$。设 $i<j$。由于 $x_i$ 紧，由引理 A.8.3 可见 $x_i^{\perp}$ 余完备且稳定；又按假设它包含 $x_j$，故它必包含 $\langle x_j\rangle=\mathcal A_j$。\ysrc{374}现设 $y\in\mathcal A_j$ 任意。我们刚刚证明了 $\operatorname{RHom}(x_i,y)=0$，故 $x_i\in{}^{\perp}y$；由引理 A.8.2，${}^{\perp}y$ 余完备且稳定。于是 ${}^{\perp}y$ 包含 $\mathcal A_i=\langle x_i\rangle$；又因 $y$ 任意，便得到对任意 $x\in\mathcal A_i$、$y\in\mathcal A_j$ 有 $\operatorname{RHom}(x,y)=0$，如所欲证。
\end{proof}

\subsubsection*{A.8.12.\quad 生成元与伴随}
现在回到引理 A.8.4 的设定。设 $F':\mathcal C\isom\mathcal B$ 是第 (1) 条的等价，$G':\mathcal B\isom\mathcal C$ 是它的拟逆。以 $\iota:\mathcal B\subseteq\mathcal D$ 记包含函子，以 $\pi:\mathcal D\to\mathcal D/\mathcal A$ 记投影。则 $F\cong\iota F'$，而 $G\cong G'(\pi\iota)^{-1}\pi$（注意由 (3)，$\pi\iota$ 是等价）。又由 (5) 与 (6)，子范畴 $\mathcal A$ 与 $\mathcal B$ 中任一个都决定另一个，而函子 $\iota$ 与 $\pi$ 由 $\mathcal A$ 与 $\mathcal B$ 典范地定义。因此，只要给出子范畴 $\mathcal A$ 或 $\mathcal B$ 之一，连同等价 $F'$，就确定了函子 $F$ 的原始数据。

设 $\{X_i\}_{i\in I}$ 是 $\mathcal C$ 的一族生成元。则
\[
\mathcal B^{\perp}=\bigcap_{i\in I}F(X_i)^{\perp}=\{Y\in\mathcal D\mid \operatorname{RHom}(F(X_i),Y)=0\ \text{对一切}\ i\in I\}.
\]
于是由引理 A.8.4(5) 可得 $\mathcal A=\bigcap_{i\in I}F(X_i)^{\perp}$。

于是，假定我们知道函子 $F$ 在生成族 $\{X_i\}_{i\in I}$ 上的取值，记 $Y_i:=F(X_i)$。则我们可以把 $\mathcal B$ 刻画为由 $Y_i$ 生成的稳定全子 $\infty$-范畴。给定 $\mathcal B$ 之后便可确定 $\mathcal A$；尽管由这些数据不能确定 $F$ 本身，却能确定复合函子 $FG$。具体地，$FG\cong\iota(\pi\iota)^{-1}\pi$。

另一种表述方式如下：$\mathcal D$ 的任意对象 $Y$ 都有半正交分解\tn{底本此句作「any object $Y$ of $\mathcal D$ semiorthogonal decomposition」，脱一动词（疑为 has），译文按文意补足。}，即存在形如
\[
X\to Y\to Z
\]
的余纤维序列（等价地，纤维序列或正合三角），其中 $X\in\mathcal B$，$Z\in\mathcal A$。该余纤维序列在等价意义下唯一：取 $X=FG(Y)$，$Z$ 为伴随 $FG(Y)\to Y$ 的余单位的余纤维即得。于是，$FG(Y)$ 正是特异三角中的项 $X$。

在 $F$ 还是猜想性的应用场合，我们可以运用这套形式。更确切地说，可以假定给定 $\mathcal D$ 的一族对象 $\{Y_i\}_{i\in I}$（我们设想它们来自 $\mathcal C$ 中一生成对象集 $\{X_i\}$ 在函子 $F$ 下的像；但实际上并不会给定 $F$）。于是我们可以像上一段那样，把 $\mathcal B$ 定义为由 $Y_i$ 生成的子范畴。按构造，包含 $\mathcal B\subset\mathcal D$ 全忠实且保持余极限，故由引理 A.4.6(1)，它有右伴随。于是引理 A.8.4 适用：可以令 $\mathcal A:=\mathcal B^{\perp}$；每个对象 $Y\in\mathcal D$ 都容许一个函子性的半正交分解，由余纤维序列
\[
X\to Y\to Z
\]
给出，其中 $X\in\mathcal B$，$Z\in\mathcal A$；而函子 $Y\mapsto X$ 是 $\mathcal B$ 的包含函子的右伴随。

\subsubsection*{A.8.13.\quad 限制到开子叠}
设进一步有一个余完备稳定 $\infty$-范畴 $\mathcal D'$，带函子 $j^*:\mathcal D\to\mathcal D'$ 以及全忠实的右伴随 $j_*:\mathcal D'\to\mathcal D$。（在应用中，$\mathcal D$ 将是叠 $\mathcal X$ 上的层范畴，$\mathcal D'$ 将是开子叠 $j:\mathcal U\to\mathcal X$ 上的层范畴，而 $j^*$、$j_*$ 取通常含义。）设 $\mathcal B'$、$\mathcal A'$ 是 $\mathcal D'$ 的稳定子 $\infty$-范畴，如上给出一个半正交分解：即 $\mathcal A'$、$\mathcal B'$ 共同生成 $\mathcal D'$，且有 $\mathcal B'={}^{\perp}\mathcal A'$。

\ysrc{375}令 $\mathcal A:=j_*\mathcal A'$。于是（由于 $j_*$ 的全忠实蕴含 $j^*j_*$ 自然等价于恒等）有 $\mathcal A'=j^*\mathcal A$；又由伴随性可得
\begin{equation}\tag{A.8.14}\label{eq:a4-18-1}
{}^{\perp}\mathcal A=(j^*)^{-1}({}^{\perp}\mathcal A')=(j^*)^{-1}(\mathcal B').
\end{equation}

\section*{附录 B. ind-代数叠上的凝聚层、ind-凝聚层与 pro-凝聚层}

\subsection*{B.1.\quad $I$-进完备环上的模}
我们将建立一些结果，把各种模的稳定 $\infty$-范畴联系起来。这些结果为形式代数叠上的（ind-）凝聚层理论提供交换代数方面的背景。事实上，我们的结果在某些非交换情形下也适用；此时我们将在下文用它们推出关于 $p$-进解析群的光滑表示的稳定 $\infty$-范畴的若干事实。

设定如下：令 $A$ 是（未必交换的）Noether 环，$I$ 是 $A$ 的双边理想，且 $A$ 关于 $I$ 是 $I$-进完备的。我们进一步假定 $I$ 满足 Artin--Rees 性质；即对任意 $n\ge 0$，$A^{\oplus n}$ 的任意 $A$-子模 $M$ 都关于 $A^{\oplus n}$ 上的 $I$-进拓扑闭，并且由 $A^{\oplus n}$ 上的 $I$-进拓扑在 $M$ 上诱导的拓扑与 $M$ 上的 $I$-进拓扑一致。由该假设可得：任意有限生成 $A$-模关于其 $I$-进拓扑都是完备且分离的；而有限生成 $A$-模之间的态射，其像关于 $I$-进拓扑闭\tn{底本此处作「have $I$-adically closed closed image」，多一 closed，疑为笔误。}，并且关于来源与目标上的 $I$-进拓扑是严格的。

若 $A$ 交换，则 Artin--Rees 性质自动成立。若 $A$ 非交换，则该性质未必成立（就我们所知），但在我们所关心的一个特殊情形中它成立。

\begin{exam}{B.1.1}\label{thm:a4-18-13}
若 $H$ 是紧 $p$-进解析群，$P$ 是 $H$ 的正规开 pro-$p$ 子群（单位元具有由这样的 $P$ 组成的邻域基），而 $I$ 是投影 $\mathcal O[[H]]\to k[H/P]$ 的核，则 $\mathcal O[[H]]$ 上、更一般地任意有限生成 $\mathcal O[[H]]$-模上的 $I$-进拓扑都与典范拓扑一致（例如见下文 §D.2 的讨论）。因此 $I$ 满足 Artin--Rees 性质。
\end{exam}

\begin{defn}{B.1.2}\label{thm:a4-18-14}
记 $D(A):=D(A\text{-}\mathrm{Mod})$；同样地，对任意 $n$ 记 $D(A/I^n):=D(A/I^n\text{-}\mathrm{Mod})$。
\end{defn}

由于 $A$ 是 Noether 环，$D(A)$ 上的 $t$-结构是凝聚的。事实上，其中心（即阿贝尔范畴 $A$-$\mathrm{Mod}$）中的紧对象恰是有限表示 $A$-模，它们生成 $A$-$\mathrm{Mod}$ 的一个阿贝尔子范畴。于是 $D(A)$ 中由凝聚对象组成的全子范畴恰是范畴 $D^b_{\mathrm{f.p.}}(A)$，它由（上同调）有界、且各阶上同调为有限表示的复形组成（因 $A$ 是 Noether 环，这等价于上同调有限生成）。把 $A$ 换成 $A/I^n$，类似的论断同样成立。

有显然的正合包含
\begin{equation}\tag{B.1.3}\label{eq:a4-18-2}
A/I^m\text{-}\mathrm{Mod}\hookrightarrow A/I^n\text{-}\mathrm{Mod}\hookrightarrow A\text{-}\mathrm{Mod},
\end{equation}
其中 $m\le n$；这些包含诱导出 $t$-正合函子
\[
D^b_{\mathrm{f.p.}}(A/I^m)\to D^b_{\mathrm{f.p.}}(A/I^n)\to D^b_{\mathrm{f.p.}}(A).
\]

\ysrc{376}由于这些函子随 $m$、$n$ 变化以显然的方式相容，我们得到一个函子
\begin{equation}\tag{B.1.4}\label{eq:a4-18-3}
\operatorname*{colim}_n D^b_{\mathrm{f.p.}}(A/I^n)\to D^b_{\mathrm{f.p.}}(A)
\end{equation}
（余极限在稳定 $\infty$-范畴构成的 $\infty$-范畴中取）。我们将在下文证明 (B.1.4) 全忠实，并描述其本质像。首先，我们考虑阿贝尔范畴层面上一个类似的构造。

\begin{defn}{B.1.5}\label{thm:a4-18-15}
以 $(A\text{-}\mathrm{Mod})_I$ 记 $A$-$\mathrm{Mod}$ 的全阿贝尔子范畴（实际上是 Serre 子范畴），它由这样的模组成：其中每个元素都被 $I$ 的某个幂零化。等价地，$(A\text{-}\mathrm{Mod})_I$ 是并 $\bigcup_n A/I^n\text{-}\mathrm{Mod}$ 在滤余极限下的闭包，其中并取在 $A$-$\mathrm{Mod}$ 中。
\end{defn}

\begin{rem}{B.1.6}\label{thm:a4-18-16}
由于 $(A\text{-}\mathrm{Mod})_I$ 是 Grothendieck 阿贝尔范畴 $A$-$\mathrm{Mod}$ 的 Serre 子范畴，且对取任意直和封闭，它本身是 Grothendieck 阿贝尔范畴。特别地，$D((A\text{-}\mathrm{Mod})_I)$ 有定义，并带有其典范 $t$-结构，其中心为 $(A\text{-}\mathrm{Mod})_I$。

范畴 $(A\text{-}\mathrm{Mod})_I$ 中的紧对象恰是有限表示模（即作为 $A$-模紧的那些模），它们生成 $(A\text{-}\mathrm{Mod})_I$ 的一个阿贝尔子范畴。于是 $D((A\text{-}\mathrm{Mod})_I)$ 上的 $t$-结构是凝聚的。其中由凝聚对象组成的全子范畴等于 $D^b_{\mathrm{f.p.}}((A\text{-}\mathrm{Mod})_I)$，即由各阶上同调为有限表示的复形组成的子范畴。
\end{rem}

还有一族范畴需要引入。

\begin{defn}{B.1.7}\label{thm:a4-18-17}
令 $D_I(A)$ 记 $D(A)$ 的全子范畴，它由各阶上同调属于 $(A\text{-}\mathrm{Mod})_I$ 的对象组成；并置 $D^b_{\mathrm{f.p.},I}(A):=D^b_{\mathrm{f.p.}}(A)\cap D_I(A)$。
\end{defn}

$D_I(A)$ 上 $t$-结构的中心恰是 $(A\text{-}\mathrm{Mod})_I$，后者局部凝聚；因此 $D^b_{\mathrm{f.p.},I}(A)$ 恰是 $D_I(A)$ 中由凝聚对象组成的全子范畴。

(B.1.3) 中的包含 $A/I^n\text{-}\mathrm{Mod}\hookrightarrow A\text{-}\mathrm{Mod}$ 显然经过 $(A\text{-}\mathrm{Mod})_I$ 分解，故 (B.1.4) 分解为
\begin{equation}\tag{B.1.8}\label{eq:a4-18-4}
\operatorname*{colim}_n D^b_{\mathrm{f.p.}}(A/I^n)\to D^b_{\mathrm{f.p.}}((A\text{-}\mathrm{Mod})_I)\to D^b_{\mathrm{f.p.},I}(A)\hookrightarrow D^b_{\mathrm{f.p.}}(A).
\end{equation}
我们将证明此分解中的前两个函子是等价。

\begin{prop}{B.1.9}\label{thm:a4-18-18}
函子
\begin{equation}\tag{B.1.10}\label{eq:a4-18-5}
\operatorname*{colim}_n D^b_{\mathrm{f.p.}}(A/I^n)\to D^b_{\mathrm{f.p.}}((A\text{-}\mathrm{Mod})_I)
\end{equation}
（即 (B.1.8) 的第一个箭头）是等价。
\end{prop}
\begin{proof}
我们将通过构造 (B.1.10) 的一个拟逆来证明。但先注意 (B.1.10) 本质满，因为在构造拟逆时要用到这一点。这种本质满性归结为如下事实：若 $X^\bullet$ 是 $(A\text{-}\mathrm{Mod})_I$ 中对象组成的有界复形，且各阶上同调都有限生成，则存在有界复形之间的拟同构 $Y^\bullet\to X^\bullet$，其中每个 $Y^n$ 都是 $(A\text{-}\mathrm{Mod})_I$ 的有限表示对象（从而 $Y^\bullet$ 实际上是某个 $n$ 的 $A/I^n$-模复形）。（这是标准的：用到 $A$ 是 Noether 环、有限生成 $A$-模的子模都有限表示；但为方便读者，我们略述证明。\ysrc{377}设 $m>n$ 时 $X^m$ 有限表示，取 $X^n$ 的一个有限生成 $A$-子模 $(X^n)'$，它同时满射到 $H^n(X^\bullet)$ 与 $\operatorname{im}(d^n:X^n\to X^{n+1})$ 上。令 $Y^\bullet$ 是 $X^\bullet$ 的子复形，其中 $r\ne n-1,n$ 时 $Y^r=X^r$，$Y^n=(X^n)'$，$Y^{n-1}=(d^{n-1})^{-1}(Y^n)$。则包含 $Y^\bullet\to X^\bullet$ 是拟同构，且对一切 $m\ge n$，$Y^m$ 有限表示；于是可对 $n$ 作降归纳。）

现在转向构造所需的拟逆。命题 A.6.12 给 $\operatorname{Ind}\operatorname*{colim}_n D^b_{\mathrm{f.p.}}(A/I^n)$ 赋予一个右完备 $t$-结构，其中心恰为 $(A\text{-}\mathrm{Mod})_I$，于是命题 A.5.7 给出一个 $t$-正合函子
\begin{equation}\tag{B.1.11}\label{eq:a4-18-6}
D^+((A\text{-}\mathrm{Mod})_I)\to\operatorname{Ind}\operatorname*{colim}_n D^b_{\mathrm{f.p.}}(A/I^n).
\end{equation}
我们断言 (B.1.11) 限制为一个函子
\begin{equation}\tag{B.1.12}\label{eq:a4-18-7}
D^b_{\mathrm{f.p.}}((A\text{-}\mathrm{Mod})_I)\to\operatorname*{colim}_n D^b_{\mathrm{f.p.}}(A/I^n),
\end{equation}
而后者还是 (B.1.10) 的拟逆。

将 (B.1.11) 与函子 $D^+(A/I^m)\to D^+((A\text{-}\mathrm{Mod})_I)$（对某个固定的 $m$）复合，得到一个 $t$-正合函子
\begin{equation}\tag{B.1.13}\label{eq:a4-18-8}
D^+(A/I^m)\to\operatorname{Ind}\operatorname*{colim}_n D^b_{\mathrm{f.p.}}(A/I^n),
\end{equation}
它在 $A/I^m$-$\mathrm{Mod}$ 上的限制与典范函子
\[
A/I^m\text{-}\mathrm{Mod}=\operatorname{Ind}D^b_{\mathrm{f.p.}}(A/I^m)^{\heartsuit}\to\operatorname{Ind}D^b_{\mathrm{f.p.}}(A/I^m)\to\operatorname{Ind}\operatorname*{colim}_n D^b_{\mathrm{f.p.}}(A/I^n)
\]
一致。典范函子
\begin{equation}\tag{B.1.14}\label{eq:a4-18-9}
D^+(A/I^m)\isom\bigl(\operatorname{Ind}D^b_{\mathrm{f.p.}}(A/I^m)\bigr)^+\to\operatorname{Ind}\operatorname*{colim}_n D^b_{\mathrm{f.p.}}(A/I^n)
\end{equation}
也有同样的性质。于是命题 A.5.7 表明，两个函子 (B.1.13) 与 (B.1.14) 本质重合。

于是，(B.1.11) 与典范函子 $D^b_{\mathrm{f.p.}}(A/I^m)\to D^+((A\text{-}\mathrm{Mod})_I)$ 的复合（后者可以分解为
\[
D^b_{\mathrm{f.p.}}(A/I^m)\to D^+(A/I^m)\xrightarrow{(B.1.13)}\operatorname{Ind}\operatorname*{colim}_n D^b_{\mathrm{f.p.}}(A/I^n)
\]
）与下述复合本质重合：
\[
D^b_{\mathrm{f.p.}}(A/I^m)\to D^+(A/I^m)\xrightarrow{(B.1.14)}\operatorname{Ind}\operatorname*{colim}_n D^b_{\mathrm{f.p.}}(A/I^n);
\]
而后者显然又与复合
\[
D^b_{\mathrm{f.p.}}(A/I^m)\to\operatorname*{colim}_n D^b_{\mathrm{f.p.}}(A/I^m)\hookrightarrow\operatorname{Ind}\operatorname*{colim}_n D^b_{\mathrm{f.p.}}(A/I^n)
\]
本质重合，其中第二个 $\operatorname{colim}_n$ 的下标照录底本\tn{底本此处作 $\operatorname{colim}_n D^b_{\mathrm{f.p.}}(A/I^m)$，被取余极限的项中指标为 $m$ 而余极限下标为 $n$，二者不一致，疑为漏排所致。}。令 $m$ 变化，并把它重新记作（哑变量）$n$，便知 (B.1.11) 与函子
\[
\operatorname*{colim}_n D^b_{\mathrm{f.p.}}(A/I^n)\xrightarrow{(B.1.10)}D^b_{\mathrm{f.p.}}((A\text{-}\mathrm{Mod})_I)\hookrightarrow D^+((A\text{-}\mathrm{Mod})_I)
\]
的复合，与包含
\[
\operatorname*{colim}_n D^b_{\mathrm{f.p.}}(A/I^n)\hookrightarrow\operatorname{Ind}\operatorname*{colim}_n D^b_{\mathrm{f.p.}}(A/I^n)
\]
在自然等价意义下重合。由于 (B.1.10) 本质满，我们还得到：(B.1.11) 在 $D^b_{\mathrm{f.p.}}((A\text{-}\mathrm{Mod})_I)$ 上的限制，其像确实落在 $\operatorname*{colim}_n D^b_{\mathrm{f.p.}}(A/I^n)$ 中，从而确实诱导出函子 (B.1.12)。此外，迄今所证表明：(B.1.12) 与 (B.1.10) 的复合与 $\operatorname*{colim}_n D^b_{\mathrm{f.p.}}(A/I^n)$ 上的恒等函子本质重合。

\ysrc{378}还需考虑 (B.1.10) 与 (B.1.12) 的复合。该复合是从 $D^b_{\mathrm{f.p.}}((A\text{-}\mathrm{Mod})_I)$ 到自身的 $t$-正合函子，它在 $D^b_{\mathrm{f.p.}}((A\text{-}\mathrm{Mod})_I)^{\heartsuit}$ 上的限制是恒等。定理 A.5.6 表明它与恒等函子本质重合。（严格说来，我们应当把定理 A.5.6 应用于 $D^+((A\text{-}\mathrm{Mod})_I)$，而非它的全子范畴 $D^b_{\mathrm{f.p.}}((A\text{-}\mathrm{Mod})_I)$。但这完全可以形式地做到：把我们的函子典范地延拓为 $\operatorname{Ind}\bigl(D^b_{\mathrm{f.p.}}((A\text{-}\mathrm{Mod})_I)\bigr)$ 的 $t$-正合自函子，再限制到 $\bigl(\operatorname{Ind}D^b_{\mathrm{f.p.}}((A\text{-}\mathrm{Mod})_I)\bigr)^+$，最后用典范等价 $D^+((A\text{-}\mathrm{Mod})_I)\isom\operatorname{Ind}\bigl(D^b_{\mathrm{f.p.}}((A\text{-}\mathrm{Mod})_I)\bigr)^+$。）
\end{proof}

在证明 (B.1.8) 中的第二个函子是等价之前，我们先回忆一个大概是标准的结果，它比较 $A$-$\mathrm{Mod}$ 与 $(A\text{-}\mathrm{Mod})_I$ 中的内射对象。

包含 $(A\text{-}\mathrm{Mod})_I\hookrightarrow A$\tn{底本此处作「$(A\text{-}\mathrm{Mod})_I\hookrightarrow A$」，按上下文与下文推论 B.1.16 的写法，应为 $(A\text{-}\mathrm{Mod})_I\hookrightarrow A\text{-}\mathrm{Mod}$。} 有右伴随，即 $M\mapsto M[I^\infty]$。由于正合函子的右伴随保持内射对象，我们看到：若 $M$ 是内射 $A$-模，则 $M[I^\infty]$ 是 $(A\text{-}\mathrm{Mod})_I$ 的内射对象。但事实上还有更强的结论，如下面的引理所示。

\begin{lem}{B.1.15}\label{thm:a4-18-19}
若 $M$ 是内射 $A$-模，则 $M[I^\infty]$ 仍是内射 $A$-模。
\end{lem}
\begin{proof}
由 Baer 判别法，只需证明：若 $J$ 是 $A$ 的理想，则任意态射 $f:J\to M[I^\infty]$ 都延拓为态射 $A\to M[I^\infty]$。由于 $J$ 有限生成（$A$ 是 Noether 环），态射 $f$ 的像落在某个 $M[I^m]$ 中，故 $f$ 经过 $J$ 的商 $J/I^mJ$ 分解。现在我们的 Artin--Rees 假设蕴含：对某个 $n\ge m$，有 $I^n\cap J\subseteq I^mJ$。于是 $f$ 经过态射 $g:J/I^n\cap J\to M[I^\infty]$ 分解。若把 $g$ 的定义域改写为 $I^n+J/I^n$\tn{底本此处印作 $I^n+J/I^n$，按文意应读作 $(I^n+J)/I^n$。}，并把它看作 $R/I^n$ 的子模\tn{底本此处以下三处把环 $A$ 印作 $R$（「$R/I^n$」「$R/I^n\to M$」「$R\to R/I^n$」），按上下文应为 $A$。}，则由 $M$ 内射的假设，$g$ 延拓为态射 $R/I^n\to M$，后者的像显然落在 $M[I^n]\subseteq M[I^\infty]$ 中。把 $g$ 的这个延拓与商映射 $R\to R/I^n$ 复合，便得到 $f$ 所需的延拓。
\end{proof}

\begin{cor}{B.1.16}\label{thm:a4-18-20}
包含 $(A\text{-}\mathrm{Mod})_I\hookrightarrow A\text{-}\mathrm{Mod}$ 保持内射对象。
\end{cor}
\begin{proof}
设 $M$ 是 $(A\text{-}\mathrm{Mod})_I$ 的内射对象，$M\hookrightarrow N$ 是把 $M$ 映入某个内射 $A$-模的包含。这个包含显然经过包含 $M\hookrightarrow N[I^\infty]$ 分解。后者是 $(A\text{-}\mathrm{Mod})_I$ 的内射对象之间的态射，故是可裂包含。另一方面，引理 B.1.15 表明 $N[I^\infty]$ 仍是内射 $A$-模，从而它的直和项 $M$ 也是。
\end{proof}

\begin{prop}{B.1.17}\label{thm:a4-18-21}
典范函子 $D^+((A\text{-}\mathrm{Mod})_I)\to D^+_I(A)$ 是等价，且限制为等价 $D^b_{\mathrm{f.p.}}((A\text{-}\mathrm{Mod})_I)\isom D^b_{\mathrm{f.p.},I}(A)$（这正是 (B.1.8) 的第二个箭头）。
\end{prop}
\begin{proof}
所断言的第一个等价由推论 B.1.16 与命题 A.7.3(1) 结合而得。把这个等价限制到各截断 $\tau^{\ge n}$（对每个 $n$）中的紧对象上，便知该等价进而限制为等价 $D^b_{\mathrm{f.p.}}((A\text{-}\mathrm{Mod})_I)\isom D^b_{\mathrm{f.p.},I}(A)$。
\end{proof}

\ysrc{379}
\subsection*{B.2.\quad ind-代数叠上的 ind-凝聚层}
本节介绍（或者说，实际上是从文献中回顾）与代数叠及形式代数叠上 $\mathcal O$-模层的导出范畴有关的一些基本概念。我们的讨论略有些不能令人满意，因为我们希望使用稳定 $\infty$-范畴，而不只是其下的三角范畴；而我们关于叠理论的基本参考文献 [Sta] 并不使用 $\infty$-范畴的概念。

例如，在 $\infty$-范畴的世界里，对于概形上拟凝聚层导出范畴（的稳定 $\infty$-范畴版本）中的对象，「等变性」是完全合用的概念；于是把代数叠 $\mathcal X$ 写成某个概形 $U$ 与群胚 $R$ 的商 $[U/R]$，就可以把 $\mathcal X$ 上拟凝聚复形的稳定 $\infty$-范畴定义为 $U$ 的相应范畴的 $R$-等变对象。（参看下文注 B.2.4。）这样一来，甚至可以在定义 $\mathcal X$ 上拟凝聚层的阿贝尔范畴之前，先定义 $\mathcal X$ 上拟凝聚复形的导出范畴；后者随后可以定义为前者上某个 $t$-结构的中心。在这个进路中，更一般的 $\mathcal O_{\mathcal X}$-模的概念，以及这种模的复形的导出范畴，根本不会出现。（诸如 $D(\mathrm{QCoh}(\mathcal X))$ 与 $D_{\mathrm{qc}}(\mathcal O_{\mathcal X})$ 之间关系这样的问题，也就不会产生。）然而在 [Sta] 的框架里，这样的定义是行不通的。

我们并不真打算在这里调和这些互相冲突的进路。我们只是陈述若干可能的定义，给出文献指引，并尽力说明它们之间的关系。归根结底，$\infty$-范畴的视角将是最适合我们目标的视角。

\subsubsection*{B.2.1.\quad 代数叠的情形}
若 $\mathcal X$ 是代数叠，则我们以 $D(\mathcal X)$ 记 $\mathcal X$ 上拟凝聚复形的稳定 $\infty$-范畴。下面简要说明其含义。（注意我们所引的文献通常在导出/三角范畴、而非稳定 $\infty$-范畴的语境中工作，但其结果容易改写到后一设定。）

定义这个范畴（对任意代数叠 $\mathcal X$）实际上（至少！）有两种可能的进路：我们可以先考虑 $\mathcal X$ 上拟凝聚层的阿贝尔范畴 $\mathrm{QCoh}(\mathcal X)$（就连这个范畴也有多种定义：用小 lisse-étale 拓扑、小 fppf 拓扑（如 [ATJLPRVG18] 中所用）、[Sta, Tag 0786] 的平坦-fppf 拓扑（这是小 fppf 拓扑的一个变体：其中考虑的是到 $\mathcal X$ 的任意平坦态射，但覆盖取为 fppf），或者大 fppf 拓扑，不过它们给出等价的范畴（例如见 [Sta, Tag 07B1]）），它是 Grothendieck 阿贝尔范畴，然后构造 $D(\mathrm{QCoh}(\mathcal X))$。

或者，我们也可以考虑更大的 Grothendieck 阿贝尔范畴 $\mathcal O_{\mathcal X}$-$\mathrm{Mod}$，即所有 $\mathcal O_{\mathcal X}$-模层（同样取 $\mathcal X$ 的某个适当拓扑），构造其相伴的稳定 $\infty$-范畴 $D(\mathcal O_{\mathcal X})$，然后（尝试）考虑由各阶上同调拟凝聚（即属于 $\mathrm{QCoh}(\mathcal X)$）的对象组成的全子范畴 $D_{\mathrm{qc}}(\mathcal O_{\mathcal X})$。若使用小拓扑，这个定义有效且合理，因为此时包含 $\mathrm{QCoh}(\mathcal X)\to\mathcal O_{\mathcal X}$-$\mathrm{Mod}$ 是正合的。若使用大拓扑，则包含 $\mathrm{QCoh}(\mathcal X)\to\mathcal O_{\mathcal X}$-$\mathrm{Mod}$ 一般不正合（见 [Sta, Tag 06VE] 第 (6) 点），因而这个构造无效；但可以改行它的一个更精细的版本 [Sta, Tag 07B5]。同样，所得范畴在典范等价意义下与构造方法无关。

\ysrc{380}有典范的 $t$-正合函子 $D(\mathrm{QCoh}(\mathcal X))\to D_{\mathrm{qc}}(\mathcal O_{\mathcal X})$。即使在 $\mathcal X$ 上合理的有限性假设之下，我们也不知道它是否为等价。不过，若 $\mathcal X$ 拟紧且几何（即具有仿射对角），则该函子确实限制为等价
\begin{equation}\tag{B.2.2}\label{eq:a4-18-10}
D^+(\mathrm{QCoh}(\mathcal X))\isom D^+_{\mathrm{qc}}(\mathcal O_{\mathcal X});
\end{equation}
这是 [ATJLPRVG18, Prop. 1.6]。

$D(\mathrm{QCoh}(\mathcal X))$ 与 $D_{\mathrm{qc}}(\mathcal O_{\mathcal X})$ 上任一 $t$-结构的中心都是范畴 $\mathrm{QCoh}(\mathcal O_{\mathcal X})$。若 $\mathcal X$ 是 Noether 代数叠，则该范畴局部凝聚，其紧对象组成的全子范畴恰是 $\mathcal X$ 上凝聚层的子范畴 $\mathrm{Coh}(\mathcal X)$。特别地，此时 $D(\mathrm{QCoh}(\mathcal X))$ 与 $D_{\mathrm{qc}}(\mathcal O_{\mathcal X})$ 上的 $t$-结构都是凝聚的。再设 $\mathcal X$ 是几何的；则等价 (B.2.2) 诱导出 $D(\mathrm{QCoh}(\mathcal X))$ 与 $D_{\mathrm{qc}}(\mathcal O_{\mathcal X})$ 中凝聚对象子范畴之间的等价。我们通过这个等价把这两个凝聚对象组成的稳定 $\infty$-范畴等同起来，并任意记之为 $D^b_{\mathrm{coh}}(\mathcal O_{\mathcal X})$。这是合理的，因为这些对象（作为凝聚 $t$-结构的凝聚对象）有界，而 $D^b_{\mathrm{coh}}(\mathcal O_{\mathcal X})$ 的中心恰是 $\mathrm{Coh}(\mathcal X)$（再次由凝聚 $t$-结构的一般性质）。

\begin{defn}{B.2.3}\label{thm:a4-18-22}
设 $\mathcal X$ 是 Noether 且几何的代数叠（于是按前述讨论，$D^b_{\mathrm{coh}}(\mathcal O_{\mathcal X})$ 无歧义地有定义）。定义 $\operatorname{Ind}\mathrm{Coh}(\mathcal X):=\operatorname{Ind}D^b_{\mathrm{coh}}(\mathcal O_{\mathcal X})$，并有时称之为 $\mathcal X$ 上 ind-凝聚复形的稳定 $\infty$-范畴。
\end{defn}

\begin{rem}{B.2.4}\label{thm:a4-18-23}
若从一开始就在 $\infty$-（而非三角）范畴的语境中工作，则还可以给出（某个版本的）$D(\mathrm{QCoh}(\mathcal X))$ 的又一种定义，即由仿射情形下降得到；例如见 [Gai, §1.1.3]。（$\infty$-范畴式的工作方式对这一种定义至关重要，因为该语境中隐含的凝聚高阶同伦，是使下降在导出层面成为有效过程所必需的。）这个稳定 $\infty$-范畴通常记作 $\mathrm{QCoh}(\mathcal X)$，事实上它在远比此一般的范围——即导出预叠——上有定义。这一视角与上述视角由 [HR17, Prop. 1.3] 调和，该文证明所得范畴等价于 $D_{\mathrm{qc}}(\mathcal O_{\mathcal X})$。

我们注意到，§5.1.7 中在刚性解析 Artin 叠 $\mathcal X$ 上给出的稳定 $\infty$-范畴 $D_{\mathrm{coh}}(\mathcal X)$ 的定义，与我们刚刚描述的 $D_{\mathrm{qc}}(\mathcal X)$ 的 $\infty$-范畴定义类似。

我们还注意到，可以用同样的风格给出 $\operatorname{Ind}\mathrm{Coh}(\mathcal X)$ 的定义，从而也在远比此一般的范围中给出定义，例如见 [Gai12]、[GR11, §2.3.1] 与 [CW24]。
\end{rem}

\subsubsection*{B.2.5.\quad ind-代数叠的情形}
设 $\mathcal X$ 是 ind-代数叠，容许一个描述
\begin{equation}\tag{B.2.6}\label{eq:a4-18-11}
\mathcal X\isom\operatorname*{colim}_n\mathcal X_n,
\end{equation}
其中各 $\mathcal X_n$ 是几何 Noether 代数叠，过渡态射 $\mathcal X_n\to\mathcal X_{n+1}$ 是闭浸入。这些闭浸入诱导出 $t$-正合函子
\[
D^b_{\mathrm{coh}}(\mathcal X_n)\to D^b_{\mathrm{coh}}(\mathcal X_{n+1}),
\]
从而（过渡到 Ind 范畴）诱导出
\[
\operatorname{Ind}\mathrm{Coh}(\mathcal X_n)\to\operatorname{Ind}\mathrm{Coh}(\mathcal X_{n+1}).
\]

\ysrc{381}\begin{defn}{B.2.7}\label{thm:a4-18-24}
定义 $\operatorname{Ind}\mathrm{Coh}(\mathcal X):=\operatorname*{colim}_n\operatorname{Ind}\mathrm{Coh}(\mathcal X_n)$，其中余极限在对象为余完备稳定 $\infty$-范畴、态射为连续函子的 $\infty$-范畴中取。
\end{defn}

\begin{rem}{B.2.8}\label{thm:a4-18-25}
定义 B.2.7 中用来定义余极限的前推函子 $\operatorname{Ind}\mathrm{Coh}(\mathcal X_n)\to\operatorname{Ind}\mathrm{Coh}(\mathcal X_{n+1})$ 是 $t$-正合的，故 $\operatorname{Ind}\mathrm{Coh}(\mathcal X)$ 自然地带有 $t$-结构；该 $t$-结构还是凝聚的，且关于它 $\operatorname{Ind}\mathrm{Coh}(\mathcal X)$ 是左正则的。

我们定义 $\mathrm{Coh}(\mathcal X)$ 为 $\operatorname{Ind}\mathrm{Coh}(\mathcal X)^{\heartsuit}$ 中紧对象组成的阿贝尔范畴。等价地，$\mathrm{Coh}(\mathcal X):=\operatorname*{colim}_n\mathrm{Coh}(\mathcal X_n)$。
\end{rem}

\begin{rem}{B.2.9}\label{thm:a4-18-26}
$\mathcal X$ 的任意两个形如 (B.2.6) 的描述都以显然的方式互相共尾，因此 $\operatorname{Ind}\mathrm{Coh}(\mathcal X)$ 与 $\mathrm{Coh}(\mathcal X)$ 在典范等价意义下与此类描述的选择无关，从而有良好定义。
\end{rem}

\begin{rem}{B.2.10}\label{thm:a4-18-27}
[Gai12] 与 [GR11, §2.3.1] 中给出的 $\operatorname{Ind}\mathrm{Coh}(\mathcal X)$ 的定义（我们已在上文注 B.2.4 中回顾过），也可以直接应用于 ind-代数叠 $\mathcal X$。我们之所以给出上面更具体的定义，只是因为我们觉得它更容易思考。
\end{rem}

\subsection*{B.3.\quad 形式代数叠上的 pro-凝聚层}
设 $\mathcal X$ 是几何 Noether 形式代数叠，通过 $\operatorname*{colim}_n\mathcal X_n\isom\mathcal X$ 被描述为 ind-代数叠，其中各 $\mathcal X_n$ 是几何 Noether 代数叠，过渡映射 $\mathcal X_n\to\mathcal X_{n+1}$ 是加厚（thickening），即诱导出下层约化子叠上同构的闭浸入。于是我们可以按定义 B.2.7 与注 B.2.8 来定义 $\operatorname{Ind}\mathrm{Coh}(\mathcal X)$ 与 $\mathrm{Coh}(\mathcal X)$。

在 $\mathcal X$ 上有一些人们可能想考虑的自然「层」，按我们的定义它们并不作为 $\mathrm{Coh}(\mathcal X)$ 的对象出现。例如，若 $A$ 是 $I$-进完备的 Noether 环，则通常的 $\operatorname{Spf}A$ 上凝聚层定义给出一个范畴，它（通过取整体截面）等价于有限生成 $A$-模的范畴；而按我们的定义，$\mathrm{Coh}(\operatorname{Spf}A)$（通过取整体截面）等价于由被某个 $I^n$ 零化的有限生成 $A$-模组成的子范畴。特别地，$\operatorname{Spf}A$ 上的结构层（对应环 $A$ 本身）不是 $\operatorname{Spf}A$ 上的凝聚层。我们转而把它等同于形式 pro-对象 “$\varprojlim_n$” $A/I^n$，并把它看作 $\operatorname{Spf}A$ 上的 pro-凝聚层。

\begin{defn}{B.3.1}\label{thm:a4-18-28}
$\mathcal X$ 上 pro-凝聚层的阿贝尔范畴定义为 $\operatorname{Pro}\mathrm{Coh}(\mathcal X)$ 的形式 pro-范畴。
\end{defn}

我们也可以考虑稳定 $\infty$-范畴 $\operatorname{Pro}D^b_{\mathrm{coh}}(\mathcal X)$，即 $\mathcal X$ 上 pro-凝聚复形的稳定 $\infty$-范畴。$D^b_{\mathrm{coh}}(\mathcal X)$ 上的通常 $t$-结构延拓为 $\operatorname{Pro}D^b_{\mathrm{coh}}(\mathcal X)$ 上的 $t$-结构，其中心（典范地等价于）$\operatorname{Pro}\mathrm{Coh}(\mathcal X)$。

\begin{exam}{B.3.2}\label{thm:a4-18-29}
若 $A$ 是 Cohen--Macaulay 的 $I$-进完备 Noether 环，$D$ 是其对偶模，则 $D$ 也给出 $\operatorname{Spf}A$ 上的一个 pro-凝聚层。更一般地，若 Noether 形式代数叠容许对偶复形，则该复形自然地可看作 $\operatorname{Pro}D^b_{\mathrm{coh}}(\widehat{\mathcal X})$ 的对象。我们在正文中采用这一观点来看待对偶层/复形，因为我们要在 pro-凝聚层的语境中应用 Grothendieck--Serre 对偶。
\end{exam}

不过我们要指出：若把对偶理论限制到 $\operatorname{Ind}\mathrm{Coh}(\mathcal X)$ 的对象上，则我们可以构造对偶复形的另一个版本，\ysrc{382}它本身是 $\operatorname{Ind}\mathrm{Coh}(\mathcal X)$ 的对象。

具体地说，若记 $\mathcal X\isom\operatorname*{colim}_n\mathcal X_n$，以 $i_{n,n+1}:\mathcal X_n\to\mathcal X_{n+1}$ 记过渡映射（假定为闭浸入），以 $i_n:\mathcal X_n\to\mathcal X$ 记典范闭浸入；再设 $\omega_n$ 记 $\mathcal X_n$ 上的对偶复形（假定存在），则自然同构 $\omega_n\isom i^!_{n,n+1}\omega_{n+1}$ 诱导出态射 $(i_{n,n+1})_*\omega_n\to\omega_{n+1}$，从而诱导出态射 $(i_n)_*\omega_n\to(i_{n+1})_*\omega_{n+1}$。若置 $\omega:=\operatorname*{colim}_n(i_n)_*\omega_n$，则 $\omega$ 是 $\operatorname{Ind}\mathrm{Coh}(\mathcal X)$ 的对象，并作为（$\operatorname{Ind}\mathrm{Coh}(\mathcal X)$ 的对象的）对偶层。

作为一个具体例子，取 $A=k[[T]]$（带其 $T$-进拓扑）。则 $k[[T]]$ 是 $A$ 的对偶模。另一方面，若写 $A=\varprojlim_n k[T]/(T^n)$，则上一段的构造给出 $A$-模 $\operatorname*{colim}_n k[T]/(T^n)$，其过渡映射由乘以 $T$ 给出；它同构于 $k((T))/k[[T]]$。这个模对应于 $\operatorname{Ind}\mathrm{Coh}(\operatorname{Spf}A)$ 的一个对象，而考虑短正合列
\[
0\to k[[T]]\to k((T))\to k((T))/k[[T]]\to0
\]
便知 $\operatorname{RHom}_A(-,k[[T]])$ 与 $\operatorname{RHom}_A(-,k((T))/k[[T]])$ 在 $\operatorname{Ind}\mathrm{Coh}(\operatorname{Spf}A)$ 的对象上（差一个平移）重合（因为 $\operatorname{RHom}_A(-,k((T)))$ 在 $T$-幂挠 $A$-模上为零）。

\section*{附录 C. 局部化与 Morita 理论}

「局部化」指的是这样的想法：把某个环 $A$ 上的模 $M$ 解释为适当类型的空间（对我们而言是概形或其变体，如代数叠、形式代数叠）上的层（某种层；对我们而言总是拟凝聚层，或其某种变体，如拟凝聚层的复形、ind-凝聚层）。基本直觉来自谱理论：若 $A$ 是 $k$-代数，则 $A$-模 $M$ 是带 $A$ 作为算子环作用的 $k$-向量空间。于是我们可以尝试按 $M$ 关于 $A$ 作用的各同时特征空间（若 $A$ 交换）来分解 $M$，或者更一般地（若 $A$ 非交换）把它分解为单 $A$-模的直和或直积分。我们可以期望单 $A$-模的集合具有某种空间的结构，并且这一分解是通过把 $M$ 铺展到单 $A$-模的空间上来实现的：相应层在给定点处的纤维，对应该单模对 $M$ 的贡献的「倍数空间」\tn{底本此处作「the multiplicity space that the measures the contribution」，多一 the，疑为笔误。}。

\subsection*{C.1.\quad 在 $\operatorname{Spec}A$ 上局部化 $A$-模（典范例子）}
代数几何中局部化最基本的例子是仿射簇上的拟凝聚层理论：若 $X=\operatorname{Spec}A$ 是仿射簇，则整体截面函子
\begin{equation}\tag{C.1.1}\label{eq:a4-18-12}
\Gamma(X,-):\mathrm{QCoh}(X)\to\Gamma(X,\mathcal O_X)\text{-}\mathrm{Mod}
\end{equation}
是等价。其拟逆当然由局部化函子 $M\mapsto\mathcal O_X\otimes_A M$ 给出。

当作用于 $\mathcal O_X$-模层时，整体截面函子的另一种描述是 $\operatorname{Hom}_{\mathcal O_X}(\mathcal O_X,-)$。于是这个局部化的基本例子说明了一个一般范式：层 $\mathcal O_X$ 自然地是 $(\mathcal O_X,A)$-双模（以某种颇为平凡的方式，即通过把 $A$ 与 $\mathcal O_X$ 的整体截面环等同），于是我们有伴随函子对 $\mathcal O_X\otimes_A-$ 与

% 本片止于印刷页 382 页末（句子中断）。下片 ch-a4-19 自印刷页 383 起，其正文以
% 「……Hom_{O_X}(O_X,-)，二者合起来诱导 A-模范畴与拟凝聚 A-模范畴之间的等价」接续本片末句。
